\clearpage
\section*{Appendix AJ. Pressure-driven retention theorem and two-dimensional findings}
\addcontentsline{toc}{section}{Appendix AJ. Pressure-driven retention theorem and two-dimensional findings}
\noindent\textbf{Companion integration.} The full nine-page typeset reproduction of the complete shorter eight-section \emph{Pressure-driven retention obstructions and exact seven-ray stress geometry} text, with the following original two vector diagrams and an independent source appendix, is supplied in the 140-page cumulative PDF. The present GitHub native source reproduces its visual findings and principal result without adding a second author/abstract block. The older 129-page archive is preserved.

\paragraph{Source-frame one-sided cone.}
For $g=15971/122900$, $H_*=22371/122900$, the sharp one-sided seven-ray $(A,C)$ projection has the six vertices
\[
\left(\pm\frac{g\rho}{2},-g\rho\right),\quad
\left(\pm\frac{(5g+H_*)\rho}{4},\frac{(H_*-g)\rho}{4}\right),\quad
\left(\pm\frac{g\rho}{2},\frac{(g+H_*)\rho}{2}\right).
\]
Thus $|A|\le51113\rho/245800$ in the \emph{source-ray basis}; under a tilted physical frame $E$ the tensor is $EHE^{\mathsf T}$ and $A$ need not equal a physical cylindrical $rz$ component.

\begin{figure}[htbp]\centering
\begin{tikzpicture}[scale=.74,>=Stealth]
\draw[->,gray!65](-4.25,0)--(4.45,0)node[right]{$A/\rho$};
\draw[->,gray!65](0,-2.95)--(0,3.15)node[above]{$C/\rho$};
\draw[thick,blue!70!black,fill=blue!8]
(-1.1696,-2.3393)--(-3.7440,.2344)--(-1.1696,2.8074)--(1.1696,2.8074)--(3.7440,.2344)--(1.1696,-2.3393)--cycle;
\draw[dashed,orange!90!black](-3.7440,.2344)--(3.7440,.2344);
\fill[orange!90!black](3.7440,.2344)circle(2.5pt);
\fill[orange!90!black](-3.7440,.2344)circle(2.5pt);
\node[font=\scriptsize] at (0,2.9) {$C=(g+H_*)\rho/2$};
\node[font=\scriptsize] at (0,-2.65) {$C=-g\rho$};
\end{tikzpicture}
\caption{Original two-dimensional sharp six-vertex coefficient slice. Axis geometry is illustrative; the exact vertices are in the preceding display. This is not a source-matched spatial stress.}
\end{figure}

\paragraph{Conditional pressure-retention obstruction.}
Let $0<h<1/6$, $p=1/2-3h$, $\tau=T-t$ and $I_\tau=[T-\tau,T-\tau/2]$. Assume Appendix AC's nontrivial inward flux, local retention and pressure-dominance hypotheses. Put
$\delta_\tau=\tau^{-p/2}\sup_{I_\tau}\|w\|_{L^2(\Omega)}\to0$ and
$\mathsf S=U\otimes w+w\otimes U+w\otimes w$. Then
\[
\sup_{I_\tau}\|\nabla\mathsf S\|_{L^2(\mathbb T^3)}
\ge c\,\delta_\tau^{-1}\tau^{p/2-1}.
\]
Under Appendix AD's additional global velocity bounds, Appendix AE strengthens this to
\[
\boxed{\sum_n\frac{\tau_n^{1/2-h}}{\delta_{\tau_n}^2}<\infty,\quad \tau_n=2^{-n}\tau_0.}
\]
Hence $\delta_\tau=O(\tau^{1/4-h/2}\sqrt{\log(1/\tau)})$ is excluded under those additional hypotheses. They have \emph{not} been shown to hold for a complete H8 source and correction.

\begin{figure}[htbp]\centering
\begin{tikzpicture}[scale=.70,>=Stealth]
\fill[blue!7](0,0)circle(2.75);
\fill[white](0,0)circle(1.90);
\fill[blue!22](0,0)circle(1.25);
\draw[thick,blue!80!black](0,0)circle(1.25);
\draw[thick,dashed](0,0)circle(1.90);
\draw[thick,orange!80!black](0,0)circle(2.75);
\foreach \a in {30,120,210,300}{\draw[very thick,->,blue!75!black]({1.55*cos(\a)},{1.55*sin(\a)})arc[start angle=\a,end angle=\a+32,radius=1.55];}
\foreach \a in {0,90,180,270}{\draw[very thick,->,orange!80!black]({2.75*cos(\a)},{2.75*sin(\a)})--({2.3*cos(\a)},{2.3*sin(\a)});}
\node[fill=white,inner sep=1pt,font=\footnotesize]at(0,0){carrier $U$};
\node[font=\scriptsize,fill=white]at(0,2.12){interface $\partial\Omega$};
\node[font=\scriptsize,fill=white]at(3.70,.9){exterior stress};
\draw[-{Stealth},orange!80!black](3.4,.7)--(2.55,.35);
\node[font=\scriptsize,fill=white]at(3.7,-1.15){nonlocal pressure};
\draw[->,orange!80!black](2.85,-.95)--(1.8,-.75);
\end{tikzpicture}
\caption{Original transverse slice geometry: the carrier vanishes near the moving interface, while nonlocal pressure can be driven by exterior quadratic velocity stress. These arrows are schematic, not a certified unforced trajectory.}
\end{figure}

\paragraph{Boundary of integration.}
The exact physical equation is
$G+w_t-L_Uw+\mathcal B(w,w)=0$ with projected $\mathcal B$.
No verified common tuple $(U,w,H,\rho,E)$ simultaneously solves this equation, the actual mixed/exterior periodic pressure, the positive compact seven-ray stress equation and terminal retention. The full unforced singularity theorem remains \textbf{OPEN}.
